% TOP_PROBLEM: {"id": 20000941, "title": "Quantifier-resolved off-diagonal semi-algebraic graph Ramsey numbers", "category_id": 2, "category": {"id": 2, "name": "combinatorics", "display_name": "Combinatorics", "description": "Counting problems, graph theory, discrete structures.", "slug": "combinatorics", "order_index": 2, "created_at": "2026-07-31T15:26:25.671Z"}, "status": "partially_solved", "problem_number": "AIM-COMBINATORICS-0066", "set_id": 14, "statusQualification": "Interprets the source n^(1+o(1)) as the intended universal upper bound; literal equality fails for constant complete-graph predicates."}
% FIRST_POSTED: 2026-10-01
% ENTRY_KIND: statement_only
% UnsolvedMath ID 20000941; source catalogue code AIM-COMBINATORICS-0066
% !TeX program = lualatex
% Definition and sources only; this file is not an LLM solution attempt.
\documentclass[11pt,a4paper]{article}
\usepackage[margin=25mm]{geometry}
\usepackage{fontspec}
\setmainfont{lmroman10-regular.otf}[BoldFont=lmroman10-bold.otf,
 ItalicFont=lmroman10-italic.otf,BoldItalicFont=lmroman10-bolditalic.otf]
\usepackage{amsmath,amssymb,mathtools}
\usepackage{unicode-math}
\setmathfont{latinmodern-math.otf}
\AtBeginDocument{\let\setminus\smallsetminus}
\usepackage{xurl,hyperref}
\hypersetup{colorlinks=true,urlcolor=blue}
\setcounter{secnumdepth}{0}
\setlength{\parindent}{0pt}
\setlength{\parskip}{5pt}
\setlength{\emergencystretch}{3em}
\begin{document}
\section*{Quantifier-resolved off-diagonal semi-algebraic graph Ramsey numbers}\label{20000941}
{\small UnsolvedMath ID 20000941 · AIM-COMBINATORICS-0066 · Combinatorics · AIM Workshop Problem Lists}
\subsection{Definitions and mathematical statement}
Fix a polynomial $f\in\mathbb R[x_1,y_1,x_2,y_2]$.
For an ordered sequence $P=(v_1,\ldots,v_N)$ of distinct points in $\mathbb R^2$, define a simple undirected graph on $[N]$ by declaring $ij$ an edge, for $i<j$, exactly when $f(v_i,v_j)\geq0$.
The order of indices removes any need to assume symmetry of $f$.
A triangle is a set of three pairwise adjacent vertices, and an independent set contains no adjacent pair.

Let $R_f(3,n)$ be the least $N$ such that every graph obtained in this way has either a triangle or an independent set of size $n$.
The off-diagonal Ramsey question asks whether, for every fixed $f$ and every $\varepsilon>0$, there exists a constant $C(f,\varepsilon)$ such that
\[
 R_f(3,n)\leq C(f,\varepsilon)n^{1+\varepsilon}
 \qquad\text{for every integer }n\geq2.
\]
Equivalently, triangle-free members of this fixed polynomial family should have independent sets of size $N^{1-o(1)}$ as $N$ grows, with the error allowed to depend on $f$.
The displayed formulation specifies the upper-bound content of the workshop shorthand $n^{1+o(1)}$.
Literal asymptotic equality cannot hold for every individual polynomial, since some polynomials define complete graphs for every configuration and therefore give a bounded threshold.
No bound uniform over polynomials of unbounded complexity is demanded.
\subsection{Short English statement}
Do triangle-free planar graphs defined by a fixed polynomial sign condition have independent sets of almost linear size?
\subsection{Sources}
\begin{itemize}
\item[S1] ULAM.AI and the UnsolvedMath Contributors, supplied problems.json, record 20000941 (AIM-COMBINATORICS-0066), AIM Workshop Problem Lists. Catalogue identity and source transcription; CC BY 4.0.
\url{https://huggingface.co/datasets/ulamai/UnsolvedMath}.
\item[S2] American Institute of Mathematics, Graph Ramsey theory, item 19. Original workshop question underlying AIM-COMBINATORICS-0066.
\url{https://www.overleaf.com/read/mnvcscjjysvg}.
\end{itemize}
\end{document}
